% Fragment to be included in Articolo_team8_working.tex.
% Suggested packages in the main preamble:
% \usepackage{amsmath,amssymb,booktabs}
%
% Rewritten on 24 October 2026 so that the text describes exactly what
% src/oos_validation.py and src/calibrate_surface.py do (see the README
% section "Corrections (24 October 2026)").  The previous version described a
% weekly design with 35 transitions, 140 forecasts, a normalized ex-ante grid
% and "no variable dated after t"; none of these matched the code.
%
% ---------------------------------------------------------------------------
% NUMBER HOOK FOR THE CORRECTED RUN (q = 0.40%, 30 origins, no warm starts).
% ../img/diagnostics_ibkr/oos_article_numbers.tex is written by
% Complete-Valuation/src/build_report_fragments.py from the corrected run
% (data/processed/team8/oos_q040_20261024): counts, R^2 against persistence,
% Diebold--Mariano statistics (Newey--West, 3 lags) and dates beating
% persistence.  oos_validation.py writes only oos_run_counts.tex and
% oos_table_rows.tex, so re-running it cannot overwrite this hook.
% The placeholders below appear only if the generated file is missing.
% ---------------------------------------------------------------------------
\IfFileExists{../img/diagnostics_ibkr/oos_article_numbers.tex}{%
  % Generated by Complete-Valuation/src/build_report_fragments.py from the
% corrected rolling run data/processed/team8/oos_q040_20261024 (q = 0.40%, equal search).
\newcommand{\OOSHookCheck}{}
\newcommand{\OOSForecastOrigins}{30}
\newcommand{\OOSCommonTargetObservations}{4667}
\newcommand{\OOSPersistenceTargetObservations}{4026}
\newcommand{\OOSRTwoPersistBS}{-423.92\%}
\newcommand{\OOSDMPersistBS}{5.51}
\newcommand{\OOSRTwoPersistH}{-59.00\%}
\newcommand{\OOSDMPersistH}{3.29}
\newcommand{\OOSRTwoPersistB}{-17.47\%}
\newcommand{\OOSDMPersistB}{1.71}
\newcommand{\OOSRTwoPersistW}{-1.84\%}
\newcommand{\OOSDMPersistW}{0.29}
\newcommand{\OOSHawkesDatesBeatPersist}{13}
}{}
\providecommand{\OOSPending}{\textbf{[pending corrected run]}}
\providecommand{\OOSHookCheck}{\textbf{[check: generated OOS numbers missing]} }
\providecommand{\OOSForecastOrigins}{\OOSPending}
\providecommand{\OOSCommonTargetObservations}{\OOSPending}
\providecommand{\OOSPersistenceTargetObservations}{\OOSPending}
\providecommand{\OOSRTwoPersistBS}{\OOSPending}
\providecommand{\OOSRTwoPersistH}{\OOSPending}
\providecommand{\OOSRTwoPersistB}{\OOSPending}
\providecommand{\OOSRTwoPersistW}{\OOSPending}
\providecommand{\OOSDMPersistBS}{\OOSPending}
\providecommand{\OOSDMPersistH}{\OOSPending}
\providecommand{\OOSDMPersistB}{\OOSPending}
\providecommand{\OOSDMPersistW}{\OOSPending}
\providecommand{\OOSHawkesDatesBeatPersist}{\OOSPending}

\section{Rolling next-dense-date conditional repricing test}
\label{sec:rolling_oos}

\subsection{What the exercise is, and what it is not}
\label{subsec:oos_nature}

The temporal validation re-prices the option cross-section observed on a
\emph{target} date with model parameters estimated on an earlier
\emph{origin} date.  Structural parameters and latent states come only from the
origin.  The target date's own spot price and its own no-look-ahead Treasury
curve, however, are used as conditioning variables when the target contracts
are priced, and the target quotes are then used for scoring.  The exercise is
therefore a \emph{conditional repricing test} of the stability of the
calibrated risk-neutral surface, given the realised underlying price and
rates; it is not an ex-ante forecast of the target surface, because variables
dated at the target enter the evaluation.  It isolates the temporal stability
of each pricing model from the separate problem of forecasting GLD.

\subsection{Admitted dates, origins and targets}
\label{subsec:oos_dates}

A date is admitted to the dense rolling exercise only if, after restricting to
calls with at least 75 calendar days to expiry, it contains at least 64 unique
valid observations and at least three expiries.  Sparse dates are excluded;
no synthetic or interpolated quotes are added.  Over 6 July--2 September 2026
this rule admits 31 dense dates, from 9 July to 2 September 2026.

Each admitted date $t_i$ ($i=1,\dots,30$) is an origin, and its target is the
\emph{next available admitted dense date} $t_{i+1}$, not the next calendar day
and not a weekly bucket.  The gap between origin and target is between one
and six calendar days (one day for most pairs, three across weekends, up to
six where intermediate dates fail the density rule).  The design yields 30
origin--target pairs per model; the last admitted date is a target only.

\subsection{Origin calibration}
\label{subsec:oos_calibration}

At every origin the four models (Black--Scholes, Heston, Bates and the exact
affine Bates--Hawkes) are calibrated to the same fixed $8\times 8$
Chebyshev--Chebyshev (CC) sample of 64 \emph{actual} quotes drawn from the
origin surface.  The objective is the mean squared price error divided by the
squared Black--Scholes vega.  Discounting uses the most recent Treasury curve
dated no later than the origin, fitted with a Nelson--Siegel--Svensson curve.

GLD carries a sponsor fee that accrues against net asset value, so its forward
is $S e^{(r-q)T}$ with $q = 0.40\%$.  Because every call price used here is
homogeneous of degree one in $(S,K)$, the carry is applied exactly at the data
layer through the forward identity
\begin{equation}
    C(S,K,T;r,q) = e^{-qT}\,C\!\left(S, K e^{qT}, T; r, 0\right),
\end{equation}
that is, strikes and prices are rescaled by $e^{qT}$ and implied volatility
and vega are recomputed on the rescaled contracts.  The same transformation is
applied to every surface loaded by the rolling exercise, so calibration,
pricing and implied volatilities are all $q$-consistent.

The calibration search policy is the following, for every origin.
\begin{itemize}
    \item Black--Scholes: bounded scalar minimisation.
    \item Heston and Bates: a differential-evolution global search followed by
    local refinement, at every origin.  Bates is in addition compared with the
    nested current-origin Heston fit (zero jumps), and the lower loss is
    retained.
    \item Bates--Hawkes: a differential-evolution search over the jump and
    Hawkes block with the diffusion parameters held at the current-origin
    Bates values, followed by full SLSQP refinements from that point and from
    the Bates-like zero-excitation point; the exact nested Bates candidate is
    retained if it is better.
\end{itemize}
No model receives a warm start from an earlier origin: prior-origin starts are
opt-in (\texttt{--warm-start}) and off in the corrected run, where Heston and
Bates were recalibrated cold after an earlier pass had warm-started them.  The
legacy run instead warm-started Bates--Hawkes only.  Each result records
whether a warm start was used and its source (\texttt{warm\_start\_used},
\texttt{initialization\_source} or \texttt{warm\_start\_source}).

\subsection{State projection and target pricing}
\label{subsec:oos_projection}

Structural parameters are frozen at the origin.  Only the latent states are
moved to the target date along their risk-neutral mean dynamics over the gap
$\Delta = (t_{i+1}-t_i)/365.25$:
\begin{equation}
    v_{t_{i+1}} = \theta + (v_{t_i}-\theta)e^{-\kappa\Delta},
    \qquad
    \lambda_{t_{i+1}} = \lambda_\infty + (\lambda_{t_i}-\lambda_\infty)
    e^{-(\beta-\alpha)\Delta},
    \quad
    \lambda_\infty = \frac{\beta\bar{\lambda}}{\beta-\alpha}.
\end{equation}
Every valid call of the \emph{full} target surface (not only 64 nodes) is
then priced with the target spot and target rates, by COS inversion with 192
terms (128 for Bates--Hawkes), and converted to Black--Scholes implied
volatility.  A pricing exception leaves that contract outside the common
support; it is counted per model in the run manifest rather than dropped
silently.

\subsection{Benchmarks and supports}
\label{subsec:oos_benchmarks}

Two benchmarks are scored on the same target contracts.
\begin{itemize}
    \item \emph{Origin mean}: the cross-sectional mean implied volatility of
    the origin surface, a single number applied to every target contract.
    \item \emph{Persistence}: the origin implied volatilities of the full
    origin surface, linearly interpolated in normalised
    (maturity, log-moneyness) coordinates and evaluated at the target
    contracts (log-moneyness relative to the target spot), strictly inside the
    convex hull of the origin points, with no extrapolation.
\end{itemize}
The two benchmarks do not use the same information as the models.  The
persistence benchmark interpolates the whole origin surface, on average about
150 quotes per date (from roughly 70 to over 600 in the legacy run), whereas
each model is calibrated on the 64 CC nodes only.  Persistence also uses the
target spot to locate target contracts in moneyness.

The \emph{common model support} of a pair is the set of target contracts with
a finite market implied volatility and finite implied volatilities under all
four models; the \emph{persistence support} is its subset inside the origin
convex hull.  In the legacy run the common support contains 4{,}667 target
observations over the 30 pairs, and the persistence support 4{,}026 of them;
comparisons with persistence are computed on the persistence support only.

\subsection{Losses and out-of-sample $R^2$}
\label{subsec:oos_metrics}

For pair $t$ and contract $j$ the error is
$e^{(m)}_{j,t} = \sigma^{\mathrm{mkt}}_{j,t} - \widehat{\sigma}^{(m)}_{j,t}$,
and the date-level loss is
$\mathrm{MSE}^{(m)}_t = n_t^{-1}\sum_{j} \left(e^{(m)}_{j,t}\right)^2$ over the
relevant support.  Mean absolute error and root mean squared error are
reported as equal-weight averages of the date-level values.  The headline
out-of-sample coefficient of determination is \emph{equal-date}, so that the
four dates with about 600 target quotes do not dominate the 26 dates with
fewer than 160:
\begin{equation}
    R^2_{\mathrm{OOS},m\mid b}
    =
    1 -
    \frac{\sum_{t} \mathrm{MSE}^{(m)}_t}{\sum_{t} \mathrm{MSE}^{(b)}_t},
\end{equation}
where $b$ is the origin-mean benchmark (on the common model support) or the
persistence benchmark (with both losses on the persistence support).  The code
also reports the pooled, observation-weighted version.  A positive value means
that the model has lower loss than the benchmark.

\subsection{Inference}
\label{subsec:oos_inference}

Predictive accuracy relative to persistence is assessed with a
Diebold--Mariano test on the date-level loss differential
$d_t = \mathrm{MSE}^{(m)}_t - \mathrm{MSE}^{(\mathrm{persist})}_t$.  The
long-run variance of $\bar d$ is estimated with a Newey--West (Bartlett)
kernel with $L=\lfloor 4(T/100)^{2/9}\rfloor$ lags ($L=3$ for $T=30$), and the
statistic $\bar d/\sqrt{\widehat{\mathrm{LRV}}/T}$ is referred to a standard
normal distribution; with only 30 dates the Harvey--Leybourne--Newbold
small-sample correction and Student-$t_{T-1}$ critical values are reported
alongside; they are computed by Complete-Valuation/src/build\_dated\_paper\_audit.py (stage \texttt{oos}).

\subsection{Results of the legacy run and status of the corrected run}
\label{subsec:oos_results}

The legacy run of September 2026 used $q=0$ and gave Bates--Hawkes alone a
warm start from earlier origins.  On 30 pairs and 4{,}667 common target
observations (4{,}026 on the persistence support) \textbf{no model beat
persistence}.  The equal-date $R^2_{\mathrm{OOS}}$ against persistence was
$-268.5\%$ for Black--Scholes, $-27.2\%$ for Heston, $-12.9\%$ for Bates and
$-3.7\%$ for Bates--Hawkes; against the origin mean it was $-2.6\%$, $66.1\%$,
$73.0\%$ and $75.8\%$.  Bates--Hawkes had a lower loss than persistence on 14
of the 30 dates.  Applying the Diebold--Mariano test above to the legacy
date-level losses (the legacy \texttt{date\_metrics.csv} in
\texttt{outputs/paper\_audit\_inputs\_20260907/}) gives $5.30$ for Black--Scholes, $2.79$ for Heston, $2.16$ for Bates and $0.68$
for Bates--Hawkes (positive: persistence more accurate; HLN-corrected $4.68$,
$2.46$, $1.91$, $0.60$).  Persistence is significantly more accurate than
Black--Scholes and Heston at the 5\% level, and than Bates under the normal
reference only; the Bates--Hawkes difference is not significant.  The legacy
exercise therefore supports, at best, that Bates--Hawkes is statistically
indistinguishable from the persistence benchmark, which itself uses more
origin information than the models.  An exploratory paired block bootstrap of
the Bates--Hawkes minus Bates daily MSE is stored in
\texttt{oos\_uncertainty.json} of the September audit
(\texttt{outputs/paper\_audit\_20260907/}).

\OOSHookCheck The corrected run ($q = 0.40\%$, the calibration policy above, pricing
failures counted) covers \OOSForecastOrigins{} origins and
\OOSCommonTargetObservations{} common target observations
(\OOSPersistenceTargetObservations{} on the persistence support).  Its
equal-date $R^2_{\mathrm{OOS}}$ against persistence is \OOSRTwoPersistBS{}
(Black--Scholes), \OOSRTwoPersistH{} (Heston), \OOSRTwoPersistB{} (Bates) and
\OOSRTwoPersistW{} (Bates--Hawkes), with Diebold--Mariano statistics
\OOSDMPersistBS{}, \OOSDMPersistH{}, \OOSDMPersistB{} and \OOSDMPersistW{};
Bates--Hawkes beats persistence on \OOSHawkesDatesBeatPersist{} dates.  The
table rows are produced by the code in
\texttt{../img/diagnostics\_ibkr/oos\_table\_rows.tex}.

\subsection{Calibration fragility}
\label{subsec:oos_fragility}

The single-date calibrations behind these comparisons are fragile, and the
parameters should not be read as precise estimates.  On the 2 September 2026
CC sample with $q=0.40\%$ the Feller condition $2\kappa\theta \ge \xi^2$,
imposed as a hard constraint, binds at the optimum for Bates and
Bates--Hawkes (gap of order $10^{-8}$) and is within 0.3\% of binding for
Heston; relaxing it lowers the Heston loss by about 68\%.  The mean-reversion
speed is weakly identified: re-optimising the other Heston parameters at fixed
$\kappa$ keeps the loss within 3\% of the optimum for $\kappa\in[2,4]$.  The
jump and Hawkes parameters are not individually identified from 64 calls;
only their combined effect on the smile is.

\subsection{Other modes and data-availability caveat}
\label{subsec:ibkr_caveat}

\texttt{oos\_validation.py} also offers a fixed-origin stress mode in which
structural parameters are frozen at one origin and only the states are
projected over each target gap; it is not part of the results above.

Interactive Brokers does not provide historical data for expired option
contracts, so a retrospective reconstruction is mechanically tilted towards
longer-dated options.  Coverage is therefore audited date by date
(\texttt{audit\_market\_data.py}) and dates that fail the density rule,
including sparse reconstructed dates, are excluded from this exercise.
